\documentclass[10pt,a4paper]{amsart}
\usepackage[margin=19mm]{geometry}
\usepackage{amsmath,amssymb,mathtools}
\usepackage[scr=rsfs,cal=euler]{mathalfa}
\usepackage{xfrac}
\usepackage[unicode,hidelinks,bookmarks=false]{hyperref}
\newcommand{\ie}{i.e.\ }
\newcommand{\cf}{cf.\ }
\newcommand{\resp}{respectively\ }
\newcommand{\Subsection}{\subsection}
\providecommand{\nobreakdash}{\nobreak\hbox{-}}
\setlength{\emergencystretch}{2em}
\multlinegap0pt
\usepackage{bm}
\usepackage{amssymb}
\usepackage{mathtools}
\usepackage{algorithm}\usepackage{algpseudocode}
\usepackage[shortlabels]{enumitem}
\usepackage{xcolor}
\usepackage{multirow}
\usepackage{booktabs}
\newtheorem{theorem}{Theorem}
\newcommand{\R}{\mathbb{R}}
\title{A Primal-Dual Level Set Method for Computing Geodesic Distances}
\author{Hailiang Liu, Laura Zinnel}
\date{Source: arXiv:2601.23244v2 (CC BY 4.0)}
\begin{document}
\maketitle

\noindent\emph{Source and licence.} A Primal-Dual Level Set Method for Computing Geodesic Distances. Version arXiv:2601.23244v2 (CC BY 4.0), distributed by arXiv under the Creative Commons Attribution 4.0 International licence, http://creativecommons.org/licenses/by/4.0/, as recorded in the arXiv licence metadata for this exact version and shown on the abstract page https://arxiv.org/abs/2601.23244v2. Full licence notice: \texttt{licenses/arxiv-2601.23244-CC-BY-4.0.txt}.\par\smallskip

\noindent\textit{Editorial scope.} Counted unit: the article's theorem theorem (source0.tex lines 1490--1513). The article's complete original proof of it is delivered with it (source0.tex lines 1514--1577, one \texttt{proof} environment of 64 lines). No mathematics is written, repaired or re-derived: the statement and the proof are the article's own lines, in the article's own notation, and no retained line is substituted or re-flowed.  No further source-local context is needed: every cross-reference of every retained line resolves inside the two delivered spans.  Nothing was omitted from any retained span, so the package carries no omission record.  No retained line contains a citation, so the wrapper carries no bibliography and no bibliography entry.  No retained line contains a figure environment, an image-inclusion command or drawing code, so no image is delivered and the package declares no asset dependency.  Editorial formatting only, in addition: an AMS \texttt{amsart} preamble that restates the article's own declarations and macros used by the retained lines, namely the environments \texttt{theorem} on separate counters, together with the standard packages those lines need.  The retained environments print in this document as Theorem 1 in order of appearance, whereas the article prints the same items with the numbering of its own full document.  The complete native source of the article as distributed in the arXiv e-print for this exact version is bundled unchanged as \texttt{sources/193695/source0.tex}.
\begin{theorem}
    Assume $\phi(\gamma) = a\cdot \gamma$ for some $a \in \R^3$, and let $\lambda_k$ and $\gamma_k$ be the iterates generated by the  scheme
    \[\begin{cases}
        \lambda_{k+1} & = \lambda_k + \tau_{\lambda}(a\cdot\gamma_k - \varepsilon\lambda_{k+1})\\
        \gamma_{k+1} &= (1-\tau_{\gamma}\partial^2)^{-1}(\gamma_k - \tau_{\gamma}(2\lambda_{k+1}-\lambda_k)a).
    \end{cases}\]
    Denote the matrix and variables 
    $$
    A = \begin{bmatrix}
    \frac{1}{\tau_{\lambda}} & a^T\\
    a & \frac{1}{\tau_{\gamma}}I_{3\times 3}
\end{bmatrix}, \quad \xi = \begin{bmatrix}
    \lambda\\ \gamma
\end{bmatrix}, \;  \bar{\xi}_k = \begin{bmatrix}
    \bar{\lambda}_k\\ \bar{\gamma}_k
\end{bmatrix} = \frac{1}{k}\sum_{i=1}^k \xi_i.
$$ 
Then the system converges for any step sizes $\tau_{\lambda}, \tau_{\gamma}$ satisfying  
\[
\tau_{\lambda}\tau_{\gamma} < \frac{1}{|a|^2} = \frac{1}{a^Ta},
\]
with an ergodic convergence rate of $\mathcal{O}(1/k)$. Specifically, for any $k\geq 1$ and any $(\gamma, \lambda)$, 
\[\mathcal{L}_{\varepsilon}(\bar{\gamma}_k, \lambda) - \mathcal{L}_{\varepsilon}(\gamma, \bar{\lambda}_k) \leq \frac{1}{2k}\|\xi_0 - \xi\|^2_A. \]
\end{theorem}

\begin{proof}
From the update scheme, we have 
\[\begin{cases}
    \frac{\lambda_k -\lambda_{k+1}}{\tau_{\lambda}}+ a\cdot(\gamma_k - \gamma_{k+1}) &= \varepsilon\lambda_{k+1}-a\cdot \gamma_{k+1} = -\frac{\delta}{\delta\lambda}\mathcal{L}_{\varepsilon}^{k+1}\\
    \frac{\gamma_k - \gamma_{k+1}}{\tau_{\gamma}} + (\lambda_k-\lambda_{k+1}) a & = -\ddot{\gamma}_{k+1}+\lambda_{k+1}a = \frac{\delta}{\delta\gamma}\mathcal{L}_{\varepsilon}^{k+1}.
\end{cases}\]
Now consider the difference in Lagrangian values:
\begin{align*}
    \mathcal{L}_{\varepsilon}(\gamma_{k+1},\lambda) - \mathcal{L}_{\varepsilon}(\gamma,\lambda_{k+1}) & = \Delta_1 + \Delta_2 + \Delta_3.
\end{align*}
Term-by-term analysis gives:
\begin{itemize}
    \item Kinetic term difference:
    \begin{align*}
    \Delta_1 & = \frac{1}{2}\int_0^1\left(|\dot{\gamma}_{k+1}|^2 -|\dot{\gamma}|^2\right)~dt\\
    & = \int_0^1 (\gamma-\gamma_{k+1})\cdot\ddot{\gamma}_{k+1}~dt - \frac{1}{2}\int_0^1|\dot{\gamma}_{k+1}-\dot{\gamma}|^2~dt,
\end{align*}
\item By-linear term difference: 
\begin{align*}
    \Delta_2 & = \int_0^1 \lambda a\cdot \gamma_{k+1}~dt -\int_0^1 \lambda_{k+1}a\cdot \gamma~dt\\
     & = \int_0^1 (\lambda - \lambda_{k+1})a\cdot \gamma_{k+1}~dt + \int_0^1\lambda_{k+1}a\cdot (\gamma_{k+1}-\gamma)~dt,
\end{align*}
\item Regularization term difference: 
\begin{align*}
    \Delta_3 & = -\frac{\varepsilon}{2}\int_0^1 \lambda^2~dt + \frac{\varepsilon}{2}\int_0^1 \lambda_{k+1}^2~dt\\
    & = \int_0^1(\lambda_{k+1} - \lambda)\varepsilon\lambda_{k+1}~dt - \frac{\varepsilon}{2}\int_0^1 (\lambda- \lambda_{k+1})^2~dt.
\end{align*}
\end{itemize}
Combining all:
\begin{align*}
    & \mathcal{L}_{\varepsilon}(\gamma_{k+1},\lambda) - \mathcal{L}_{\varepsilon}(\gamma,\lambda_{k+1}) \\
    \qquad  = &\int_0^1 (\gamma - \gamma_{k+1})\cdot (\ddot{\gamma}_{k+1} - \lambda_{k+1}a)~dt +\int_0^1 (\lambda-\lambda_{k+1})(a\cdot \gamma_{k+1} - \varepsilon\lambda_{k+1})~dt \\
      \qquad  &- \frac{1}{2}\int_0^1 |\dot{\gamma} - \dot{\gamma}_{k+1}|^2~dt - \frac{\varepsilon}{2}\int_0^1 |\lambda - \lambda_{k+1}|^2~dt.
\end{align*}
Using the update expressions, 
\begin{align*}
     \leq & \int_0^1 (\gamma_{k+1}-\gamma)\cdot \left[\frac{\gamma_k -\gamma_{k+1}}{\tau_{\gamma}} + a(\lambda_k -\lambda_{k+1})\right]~dt \\
    &+ \int_0^1(\lambda_{k+1}- \lambda)\left[\frac{\lambda_k - \lambda_{k+1}}{\tau_{\lambda}}+ a\cdot(\gamma_k - \gamma_{k+1})\right]~dt.
  \end{align*}   
  Let $A = \begin{bmatrix}
    \frac{1}{\tau_{\lambda}} & a^T\\
    a & \frac{1}{\tau_{\gamma}}I_{3\times 3}
\end{bmatrix}$, and define $\xi = \begin{bmatrix}
    \lambda\\ \gamma
\end{bmatrix}$, then we obtain 
  \begin{align*} 
     =& \int_0^1 A(\xi_k - \xi_{k+1})\cdot (\xi_{k+1}-\xi)~dt \\
     = & \frac{1}{2}\int_0^1 A(\xi_k - \xi)\cdot (\xi_k - \xi)~dt - \frac{1}{2}\int_0^1A(\xi_{k+1} - \xi)\cdot (\xi_{k+1}- \xi)~dt\\ &- \frac{1}{2}\int_0^1 A(\xi_k - \xi_{k+1})\cdot (\xi_k- \xi_{k+1})~dt\\
     \leq &\frac{1}{2}\int_0^1 \|\xi_k - \xi\|_A^2~dt - \frac{1}{2}\int_0^1 \|\xi_{k+1}- \xi\|_A^2~dt.
\end{align*}
Hence, the algorithm converges if
\[\tau_{\lambda}\tau_{\gamma} < \frac{1}{|a|^2} = \frac{1}{a^Ta}.\] 
Further, define ergodic averages,    $\bar{\lambda}_k = \frac{1}{k}\sum_{i=1}^k \lambda_i$ and 
$\bar{\gamma}_k = \frac{1}{k}\sum_{i=1}^k \gamma_i$. 
By Jensen's inequality: 
\[\int_0^1 |\dot{\bar{\gamma}}_k|~dt  \leq \frac{1}{k}\sum_{i=1}^k\int_0^1 |\dot{\gamma_i}|~dt.\]
Using the telescoping sum, we conclude  
\begin{align*}
    \mathcal{L}_{\varepsilon}(\bar{\gamma}_k, \lambda) - \mathcal{L}_{\varepsilon}(\gamma, \bar{\lambda}_k) 
    & \leq \frac{1}{k}\sum_{i=0}^{k-1} (\mathcal{L}(\gamma_{i+1}, \lambda) -  \mathcal{L}(\gamma,\lambda_{i+1}))\\
     & \leq \frac{1}{2k}\|\xi_0 -  \xi\|^2_A. 
\end{align*}
Thus, the method achieves ergodic convergence at rate $O(1/k)$.
\end{proof}

\end{document}
